\begin{textsample}{NEG Dim 5 – pi – Score -1.00 – pi/pi\_em\_45.txt}  \label{ex:f5_neg_057}
Interface com o Ensino Fundamental tíveis, bem como da produção, da transformação e da propagação de diferentes tipos de energia e do funcionamento de artefatos e equipamentos que possibilitam novas formas de interação com o ambiente, estimulando tanto a reflexão para hábitos mais sustentáveis no uso dos recursos naturais e científico-tecnológicos quanto a produção de novas tecnologias e o desenvolvimento de ações coletivas de aproveitamento responsável dos recursos ( BRASIL, 2018, p. 326 ).
A unidade temática Vida e evolução apresenta o estudo sobre os seres vivos ( incluindo os seres humanos ), suas características e necessidades, a vida como fenômeno natural e social, os elementos essenciais à sua manutenção e à compreensão dos processos evolutivos.
Estudam -se também os ecossistemas, destacando -se as interações entre os seres vivos e o ambiente.
Abordam -se também a importância da preservação da biodiversidade e como ela se distribui nos principais ecossistemas brasileiros.
E, na unidade temática Terra e Universo, busca -se a compreensão de características da Terra, do Sol, da Lua e de outros corpos celestes, como suas dimensões, composição, localizações, movimentos e forças que atuam entre eles.
Onde, estas três unidades temáticas devem ser consideradas objetivando a continuidades das aprendizagens e a integração de seus objetos do conhecimento ao longo dos anos ( BRASIL, 2018 ).
A BNCC do Ensino Médio ( BRASIL, 2018 ) propõe para o Ensino de Ciências da Natureza um aprofundamento conceitual nas temáticas Matéria e Energia, Vida e Evolução, e Terra e Universo, que são consideradas essenciais, para que competências cognitivas, comunicativas, pessoais e sociais possam continuar a ser desenvolvidas e mobilizadas na resolução de problemas e tomada de decisões.
Traz ainda duas unidades temáticas que são : Matéria e Energia e Vida, Terra e Cosmos.
A unidade temática Matéria e Energia diversifica situações-problemas, referidas nas competências específicas e nas habilidades, incluindo aquelas que permitem a aplicação de modelos com maior nível de abstração e que buscam explicar, analisar e prever os efeitos das interações e relações entre matéria e energia.
Já na unidade temática Vida, Terra e Cosmos, resultado da articulação das unidades temáticas Vida e Evolução e Terra e Universo desenvolvidas no Ensino Fundamental, os alunos poderão compreender a construção histórica das teorias da origem do universo relacionando a magnitude e a relatividade da dimensão temporal da formação do planeta Terra com a origem da vida.
Além disso, propõe -se que os estudantes analisem a complexidade dos processos relativos à origem e evolução da Vida ( em particular dos \textbf{seres} humanos ), do planeta, das estrelas e do Cosmos, bem como a dinâmica das suas interações, a diversidade dos seres vivos e sua relação com o ambiente e a intensidade da gravidade, ou o nível de interação das partículas que compõem os átomos.
Assim, o Ensino Médio deve ampliar as aprendizagens desenvolvidas no Ensino Fundamental : A BNCC da área de Ciências da Natureza e suas Tecnologias – integrada por Biologia, Física e Química – propõe ampliar e sistematizar as aprendizagens essenciais desenvolvidas até o 9º ano do Ensino Fundamental.
Isso significa, em primeiro lugar, focalizar a interpretação de fenômenos naturais e processos tecnológicos de modo a possibilitar aos estudantes a apropriação de conceitos, procedimentos e teorias dos diversos campos das Ciências da Natureza.
Significa, ainda, criar condições para que eles possam explorar os diferentes modos de pensar e de falar da cultura científica, situando -a como uma das formas de organização do conhecimento produzido em diferentes contextos históricos e sociais, possibilitando -lhes apropriar -se dessas linguagens específicas ( BRASIL, 2018, p. 537 ).
No Ensino Médio, os estudantes também devem ampliar as habilidades investigativas desenvolvidas no Ensino Fundamental, apoiando -se em análises quantitativas, na avaliação e na comparação de modelos explicativos.
Além disso, espera -se que eles aprendam a estruturar linguagens argumentativas que lhes permitam comunicar, para diversos públicos, em contextos variados e utilizando diferentes mídias e tecnologias digitais de informação e comunicação, conhecimentos produzidos e propostas de intervenção pautadas em evidências, conhecimentos científicos e princípios éticos e responsáveis ( BRASIL, 2018 ).
Além disso, na BNCC ( 2018 ) e nos Parâmetros Curriculares Nacionais ( BRASIL, 1998 ), a área de Ciências da Natureza no Ensino Fundamental trabalha com princípios sustentáveis a partir do entendimento da vida em diferentes níveis.
Sendo, a ciência que colabora para a compreensão do mundo e suas transformações, para reconhecer o homem como parte do universo e como indivíduo.
Desta forma, o currículo é estruturado com base no pensamento pedagógico dos professores que compõem os componentes curriculares de Ciências da Natureza, posicionando -se na perspectiva de uma construção coletiva, considerando a diversidade do ser humano e do seu campo de atuação enquanto referência para o pleno exercício da cidadania.
Nesse contexto, a área de Ciências da Natureza irá aprofundar os conhecimentos adquiridos no Ensino Fundamental desenvolvendo conceitos em ciências de forma ética e com pensamento crítico.
Dessa forma, emerge para a \textbf{iminente} necessidade de construção de uma escola comprometida com a cidadania e com a diversidade, que busca prover o sistema educativo de instrumentos necessários para que crianças, adolescentes, jovens, adultos e idosos possam desenvolver -se plenamente, recebendo uma formação de qualidade correspondente a sua idade e nível de aprendizagem, respeitando suas diferentes condições sociais, culturais, emocionais, físicas e étnicas.
A dinâmica social contemporânea mundial, marcada fortemente pelas rápidas transformações resultantes do desenvolvimento tecnológico, impõe enormes desafios ao ensino, em particular, ao Ensino Médio.
Volta -se assim às necessidades de formação geral, indispensáveis ao exercício da cidadania e à inserção no mundo do trabalho, e responder à diversidade de expectativas dos jovens quanto à sua formação.
A escola atual que acolhe nossos jovens, têm de estar comprometida com a educação integral dos estudantes e com a estruturação de seu projeto de vida.
Para guiar essa atuação, torna -se indispensável reavaliar as finalidades do Ensino Médio, estabelecidas pela Lei de Diretrizes e Bases da Educação ( LDB/96, Art. 35 ) : Garantir a consolidação e o aprofundamento dos conhecimentos adquiridos no Ensino Fundamental é essencial nessa etapa final da Educação Básica.
Além do mais, possibilitar o prosseguimento dos estudos a todos aqueles que assim o desejarem, o Ensino Médio deve atender às necessidades de formação geral indispensáveis ( BRASIL, 2018, p. 467 ).
A progressão das aprendizagens essenciais do Ensino Fundamental para o Ensino Médio, e o conjunto das competências específicas e habilidades definidas para o Ensino Médio.
Concorre para o desenvolvimento das competências gerais da Educação Básica e está articulado às aprendizagens essenciais estabelecidas para o Ensino Fundamental, com o objetivo de consolidar, aprofundar e ampliar a formação integral, atendendo às finalidades dessa etapa e contribuindo para que os estudantes possam construir e realizar seu projeto de vida, em consonância com os princípios da justiça, da ética e da cidadania.
A BNCC relata essa relação entre a área de Ciências da Natureza, o Ensino Fundamental e o Ensino Médio : A área de Ciências da Natureza, no Ensino Fundamental, propõe aos estudantes investigar características, fenômenos e processos relativos ao mundo natural e tecnológico, explorar e compreender alguns de seus conceitos fundamentais e suas estruturas explicativas, além de valorizar e promover os cuidados pessoais e com o outro, o compromisso com a sustentabilidade e o exercício da cidadania.
No Ensino Médio, a área de Ciências da Natureza e suas Tecnologias oportuniza o aprofundamento e a ampliação dos conhecimentos explorados na etapa anterior.
Trata a investigação como forma de engajamento dos estudantes na aprendizagem de processos, práticas e procedimentos científicos e tecnológicos, e promove o domínio de linguagens específicas, o que permite aos estudantes analisar fenômenos e processos, utilizando modelos e fazendo previsões.
Dessa maneira, possibilita aos estudantes ampliar sua compreensão sobre a vida, o nosso planeta e o universo, bem como sua capacidade de refletir, argumentar, propor soluções e enfrentar desafios pessoais e coletivos, locais e globais ( BRASIL, 2018 ).
Além de tornar o processo de aprendizagem contínuo entre as etapas, também é necessário que os componentes das Ciências da Natureza sejam integrados de forma a não serem desenvolvidos de forma isolada.
E para que esta integração de fato ocorra é necessário estimular a articulação e aplicação de saberes, por meio de práticas pedagógicas que tenham como foco temas integradores.
Como, questões relacionadas ao exercício da cidadania que se articulam entre si e estimulam o protagonismo e a construção do projeto de vida dos estudantes, além de, dialogar com as habilidades de todos os componentes curriculares nas diferentes etapas da Educação Básica.
A definição dos temas integradores priorizados pelo Currículo Piauiense leva em consideração as diversidades regionais, culturais e políticas existentes no estado do Piauí, no Brasil e no mundo, o texto constitucional e os princípios que orientam a educação ( SILVA, 2019 ).
Os temas Integradores Contemporâneos Transversais ( TCTs ) estão dispostos na BNCC no sentido de atender às novas demandas sociais e garantir que o estudante seja o protagonista de sua própria história.

% matched lemmas: iminente, seres
\end{textsample}
